%!TeX root = BasiDiDati.tex

\chapter{Tecnologie per le basi di dati}

\definizione{DBMS}{
  Un \textbf{D}atabase \textbf{M}anagement \textbf{S}ystem è un software in grado di gestire collezioni di dati che siano \textcolor{azzurrino}{grandi}, \textcolor{azzurrino}{condivise} e \textcolor{azzurrino}{persistenti}, assicurando la loro \textcolor{azzurrino}{affidabilità} e \textcolor{azzurrino}{privatezza}.

  \noindent\rule{\textwidth}{0.4pt}

  \medskip
  \begin{itemize}[label=$\triangleright$, itemsep=\smallskipamount]
    \item \textbf{grandi}: dati devono essere organizzati in modo efficiente per poter essere gestiti in grandi quantità
    \item \textbf{condivise}: dati devono essere accessibili da più utenti contemporaneamente
    \item \textbf{persistenti}: dati devono essere memorizzati in modo permanente, anche dopo la chiusura del programma
    \item \textbf{affidabili}: dati devono essere protetti da errori e guasti del sistema
    \item \textbf{privati}: dati devono essere protetti da accessi non autorizzati 
  \end{itemize}
}

Un DBMS in quanto sistema software deve essere \textbf{efficiente} ed \textbf{efficace} e deve, inoltre, estendere le funzionalità del file system.


\definizione{base di dati}{Una \textbf{base di dati} è una collezione di dati gestita da un DBMS.}

\section{Architettura transazionale di un DBMS}
Un DBMS è costituto da tra 3 livelli:

\begin{figure}[H]
  \centering
  \includegraphics[width=0.8\textwidth]{images/schema.png}
\end{figure}

Un DBMS basato sul \textbf{modello relazionale} è nella maggior parte dei casi anche un sistema transazionale: fornisce un meccanismo per la definizione ed esecuzione di transazioni.

\begin{figure}[H]
  \centering
  \includegraphics[width=0.7\textwidth]{images/DBMS.png}
  \caption{sistema transazionale}
\end{figure}

\section{Transazione}

\definizione{Transazione}{Una \textbf{transazione} è un’unità di lavoro svolta da un programma applicativo (che interagisce con una base di dati) per la quale si vogliono garantire proprietà di correttezza, robustezza e isolamento.}

Le caratteristiche principali sono:
\begin{itemize}[label=$\triangleright$, itemsep=\smallskipamount]
  \item \textbf{atomicità}: una transazione  deve essere eseguita completamente o non essere eseguita affatto
  \item \textbf{non sono ammesse esecuzioni parziali}
\end{itemize}

\subsection{Sintassi}
La sintassi per definire una transazione in SQL:
\begin{codiceSQL}
<transazione> -> begin transaction
                  <programma>
                 end transaction

<programma> -> < <istruzione> | commit work | rollback work >
               { < <istruzione> | commit work | rollback work > }
\end{codiceSQL}

La transazione va a buon fine all’esecuzione di un $\verb|commit work|$, ma non ha alcun effetto se viene eseguito un $\verb|rollback work|$.

\medskip
Per essere ben formata, una transazione deve:
\begin{itemize}[label=$\circ$, itemsep=\smallskipamount]
  \item \textbf{iniziare} con un $\verb|begin transaction|$
  \item \textbf{terminare} con un $\verb|end transaction|$
  \item \textbf{esecuzione} comportare il raggiungimento di un $\verb|commit|$ o di un $\verb|rollback work|$, a seguito del quale non si eseguono altri accessi alla base di dati.
\end{itemize}

\begin{codiceSQL}
begin transaction;
  update CONTO set saldo = saldo - 1200
    where filiale = '005' and numero = 15;
  update CONTO set saldo = saldo + 1200
    where filiale = '005' and numero = 205;
  commit work;
end transaction;
\end{codiceSQL}

\subsection{Proprietà delle transazioni}
Una transazione ha 4 proprietà fondamentali:
\begin{enumerate}[itemsep=\smallskipamount]
  \item \textcolor{azzurrino}{\textbf{atomicità}}: una transazione è un'unità di esecuzione indivisibile (viene eseguita completamente o non viene eseguita affatto).
  
  Questo implica:
  \begin{itemize}[label=$\diamond$, itemsep=\smallskipamount]
    \item se una transazione viene interrotta prima del $\verb|commit|$, il lavoro fin qui eseguito deve essere disfatto ripristinando la situazione in cui si trovava la base di dati prima dell’inizio della transazione
    \item se una transazione viene interrotta dopo l’esecuzione del $\verb|commit|$ ($\verb|commit|$ eseguito con successo), il sistema deve assicurare che la transazione abbia effetto sulla base di dati
  \end{itemize}
  \item \textcolor{azzurrino}{\textbf{consistenza}}: esecuzione di una transazione non deve violare i vincoli di integrità.
  
  Questo implica:
  \begin{itemize}[label=$\diamond$, itemsep=\smallskipamount]
    \item verifica immediata:
    \begin{itemize}[itemsep=\smallskipamount]
      \item fatta nel corso della transazione
      \item viene abortita solo l’ultima operazione e il sistema restituisce all’applicazione una segnalazione d’errore
      \item applicazione può reagire alla violazione
    \end{itemize}
    \item verifica differita:
    \begin{itemize}[itemsep=\smallskipamount]
      \item  al $\verb|commit|$, se un vincolo di integrità viene violato, la transazione viene abortita senza possibilità da parte dell’applicazione di reagire alla violazione
    \end{itemize}
  \end{itemize}
  \item \textcolor{azzurrino}{\textbf{isolamento}}: esecuzione di una transazione deve essere indipendente dalla contemporanea esecuzione di altre transazioni
  \begin{itemize}[label=$\diamond$, itemsep=\smallskipamount]
    \item $\verb|rollback|$ di una transazione non deve creare $\verb|rollback|$ a catena di altre transazioni che si trovano in esecuzione contemporaneamente
    \item  sistema deve regolare l’esecuzione concorrente con meccanismi di controllo dell’accesso alle risorse
  \end{itemize}
  \item \textcolor{azzurrino}{\textbf{persistenza}}: effetto di una transazione che ha eseguito il $\verb|commit|$ non deve andare perso
  \begin{itemize}[label=$\diamond$, itemsep=\smallskipamount]
    \item sistema deve essere in grado, in caso di guasto, di garantire gli effetti delle transazioni che al momento del guasto avevano già eseguito un $\verb|commit|$
  \end{itemize}
\end{enumerate}

\smallskip
Un DBMS che gestisce transazioni dovrebbe garantire per ogni transazione che esegue tutte queste proprietà.

\section{Architettura di un DBMS}
L’architettura di un DBMS è composta da \textbf{moduli} che si occupano di gestire le diverse funzionalità del sistema.

\begin{figure}[H]
  \centering
  \includegraphics[width=0.7\textwidth]{images/ArchitetturaDBMS.png} 
  \caption{architettura di riferimento di un DBMS}
\end{figure}

Il modulo di gestione delle transazioni si occupa di trasformare una transazione in un insieme di operazioni sulla base di dati, garantendo le proprietà ACID.

\begin{figure}[H]
  \centering
  \includegraphics[width=0.7\textwidth]{images/Architettura2.png} 
  \caption{modulo di gestione delle transazioni di un DBMS}
\end{figure}

I singoli moduli che compongono questo modulo sono:
\begin{itemize}[label=$\triangleright$, itemsep=\smallskipamount]
  \item \textcolor{azzurrino}{\textbf{gestore delle interrogazioni}}: riceve in ingresso una query SQL
  \begin{figure}[H]
    \centering
    \includegraphics[width=0.7\textwidth]{images/passaggio1.png} 
  \end{figure}
  \item\textcolor{azzurrino}{ \textbf{gestore dei metodi di accesso}}: determina quali sono le pagine della memoria secondaria di cui si ha bisogno, ma capisce anche se è possibile accederci
  \begin{figure}[H]
    \centering
    \includegraphics[width=0.7\textwidth]{images/passaggio2.png} 
  \end{figure}
  \item \textcolor{azzurrino}{\textbf{gestore dell’affidabilità}}: carica in memoria principale le pagine della memoria secondaria di cui si ha bisogno e le mantiene in memoria finché non sono più necessarie
  \begin{figure}[H]
    \centering
    \includegraphics[width=0.7\textwidth]{images/passaggio4.png} 
  \end{figure}
  \item \textcolor{azzurrino}{\textbf{gestore del buffer}}
  \item \textcolor{azzurrino}{\textbf{gestore dell’esecuzione concorrente}}
  \item \textcolor{azzurrino}{\textbf{gestore dell’affidabilità}}
\end{itemize}

\bigskip
I moduli che contribuiscono a garantire le proprietà ACID sono:
\begin{itemize}[label=$\circ$, itemsep=\smallskipamount]
  \item \textbf{gestore dei metodi d'accesso}: consistenza
  \item \textbf{gestore dell'esecuzione concorrente}: atomicità e isolamento
  \item \textbf{gestore dell'affidabilità}: atomicità e persistenza  
\end{itemize}